%% new-paper.tex — START HERE
%% ---------------------------------------------------------------------------
%% A skeleton reading paper. Replace the questions with your own, then:
%%
%%     tectonic -X compile new-paper.tex
%%
%% Rules to remember:
%%   • Never type a question number. They number themselves, in source order.
%%   • The [square brackets] after each macro name are the ANSWER. It is never
%%     printed — it goes into new-paper.ans for the markscheme.
%%   • Set \ibsetmaxmarks to the total you intend. If the questions don't add
%%     up, the build warns you at the end.
%%   • Matching blocks and true-statement questions need exactly twice as many
%%     options as items/marks. The build warns you if they don't.
%%
%% Every question type is in specimens.tex if you need one that isn't here.
%% ---------------------------------------------------------------------------
\documentclass[11pt,a4paper]{article}
\usepackage{ibenglishb}          % add [answers] to print the answers in red

\ibsetsession{QDHS\,--\,2026}    % code printed top-right on every page
\ibsetmaxmarks{20}               % the skeleton below totals 20 marks.
                                 % Change to 40 (HL reading) or 25 (HL
                                 % listening) once the paper is written;
                                 % `make check` tells you when it matches.

\begin{document}

% ═══════════════════════════════════════ COVER ═════════════════════════════
\ibcover
  {English B -- Higher level -- Paper 2 -- Reading comprehension}
  {Date goes here}
  {1 h}
  {\begin{ibinstructions}{Question and answer booklet -- Instructions to students}
     \item Write your session number in the boxes above.
     \item Do not open this examination paper until instructed to do so.
     \item Answer all questions. Each question is allocated \textbf{[1 mark]}
           unless otherwise stated.
     \item Answers must be written within the answer boxes provided.
     \item All answers must be based on the appropriate texts in the
           accompanying text booklet.
     \item The maximum mark for this examination paper is \textbf{[40 marks]}.
   \end{ibinstructions}}
  {7 pages}
  {\ibsessioncode}

% ═══════════════════════════════════════ TEXT A ════════════════════════════
\ibtextheading{A}{Headline of your first text}

% --- short answer ----------------------------------------------------------
\ibrubricshort
\ibshortq[the answer]{Your question here?}
\ibshortq[the answer]{Your question here?}

% --- complete the sentence from the text -----------------------------------
% \ibparas{3}{4} prints "paragraphs [3]–[4]"; use \iblines{12}{23} instead if
% the text booklet is numbered by line.
\ibrubriccomplete{\ibparas{3}{4}}
\ibcompleteq[quoted words from the text]{A paraphrased sentence opening\dots}

% --- multiple choice (four options for reading, three for listening) --------
\ibrubricmcq
\begin{ibmcq}[C]{Your question here?}
  \opt{First option}
  \opt{Second option}
  \opt{Third option}
  \opt{Fourth option}
\end{ibmcq}

\ibturnover        % prints "Turn over" at the foot of the page being finished
\clearpage

% ═══════════════════════════════════════ TEXT B ════════════════════════════
\ibtextheading{B}{Headline of your second text}

% --- find the word ---------------------------------------------------------
\ibrubricfindword{\iblines{1}{15}}
\ibfindwordq[word from the text]{definition}
\ibfindwordq[word from the text]{definition}

% --- matching headings: \mgap numbers itself and prints [ – n – ] -----------
% Put the same [ – n – ] markers in the text booklet where the headings were.
\ibrubricheadings
\begin{ibmatch}
  \mgap[E] \mgap[B] \mgap[G] \mgap[C]
  \mopt{Heading A} \mopt{Heading B} \mopt{Heading C} \mopt{Heading D}
  \mopt{Heading E} \mopt{Heading F} \mopt{Heading G} \mopt{Heading H}
\end{ibmatch}

% --- matching vocabulary: \mitem carries its own left-hand text -------------
\ibrubricvocab
\begin{ibmatch}
  \mitem[C]{word \ibline{31}}
  \mitem[E]{word \ibline{38}}
  \mitem[A]{word \ibline{41}}
  \mopt{gloss A} \mopt{gloss B} \mopt{gloss C}
  \mopt{gloss D} \mopt{gloss E} \mopt{gloss F}
\end{ibmatch}

\clearpage

% ═══════════════════════════════════════ TEXT C ════════════════════════════
\ibtextheading{C}{An extract from \textit{Your literary text}}

% --- identify true statements: n marks, 2n options, one question number -----
\begin{ibtrue}[A, B, D, H]{4}
  \opt{Statement A} \opt{Statement B} \opt{Statement C} \opt{Statement D}
  \opt{Statement E} \opt{Statement F} \opt{Statement G} \opt{Statement H}
\end{ibtrue}

% --- true / false with justification ---------------------------------------
\ibrubrictf
\ibtfq[T: ``quotation from the text'']{Your statement here.}
\ibtfq[F: ``quotation from the text'']{Your statement here.}

% --- what the underlined words refer to ------------------------------------
\ibrubricrefer
\ibreferq[the referent]{\dots quoted phrase with \ibu{the word} underlined\dots\ \ibline{15}}

\ibblankpage      % "Please do not write on this page."

\end{document}
